% ===========================================================================
% Topology-Aware Crack and Spalling Segmentation in Concrete Dams
% via Dynamic Snake Convolution and Skeleton Recall Loss
% ===========================================================================
\documentclass[10pt,twocolumn,letterpaper]{article}

\usepackage{times}
\usepackage{epsfig}
\usepackage{graphicx}
\usepackage{amsmath}
\usepackage{amssymb}
\usepackage{booktabs}
\usepackage{multirow}
\usepackage{xcolor}
\usepackage{hyperref}
\usepackage{enumitem}

% Page geometry (CVPR-like)
\usepackage[margin=1in]{geometry}

\begin{document}

\title{Topology-Aware Crack and Spalling Segmentation in Concrete Dams\\via Dynamic Snake Convolution and Skeleton Recall Loss}

\author{
Anonymous Author(s)
}

\date{}
\maketitle

% ===========================================================================
\begin{abstract}
Automated segmentation of cracks and spalling on concrete dam surfaces is essential for structural health monitoring, yet existing methods optimize pixel-level metrics without explicitly enforcing the topological connectivity critical for safety assessment.
We propose \textbf{DSCFormer}, a topology-aware segmentation framework that augments SegFormer-B2 with two complementary enhancements: (1)~a \emph{Dynamic Snake Convolution} (DSConv) branch that performs deformable sampling along crack directions to capture elongated geometry, and (2)~\emph{Skeleton Recall Loss} (SRL) that penalizes predictions with low probability at ground-truth skeleton pixels, directly enforcing topological connectivity using precomputed skeletons with no runtime overhead.
On our \textbf{DamSegment} benchmark (1,500 images, 3~classes, 3~difficulty tiers), DSCFormer achieves 71.4\% mIoU$_{fg}$ (72.3\% with test-time augmentation), 72.1\% Boundary F1 (+5.3 over baseline), and 86.2\% clDice (+4.8 over baseline), outperforming U-Net, DeepLabV3+, Mask2Former, SegFormer-B2, and parameter-efficient fine-tuned foundation models (SAM-LoRA, DINOv2-LoRA) on all topology-sensitive metrics.
Ablations confirm that DSConv and SRL provide complementary gains: DSConv primarily improves crack connectivity (ConnR~+3.8) through geometry-aware sampling, while SRL adds consistent improvements across all metrics (mIoU~+1.0, BF1~+2.6).
Our comparison with foundation models further demonstrates that task-specific architectural design outperforms generic adaptation on small-scale industrial inspection datasets.
\end{abstract}

% ===========================================================================
% INTRODUCTION AND RELATED WORK (from Literature Review Agent)
% ===========================================================================
% ===========================================================================
% Introduction and Related Work
% Paper: Topology-Aware Crack and Spalling Segmentation in Concrete Dams
%        via Dynamic Snake Convolution and Skeleton Recall Loss
% ===========================================================================

\section{Introduction}
\label{sec:intro}

Concrete dams are safety-critical infrastructure whose structural integrity must be continuously monitored throughout multi-decade service lives.
Surface defects---predominantly cracks and spalling---serve as early indicators of internal deterioration, yet their manual inspection remains labour-intensive, subjective, and logistically constrained by the scale and accessibility of dam faces~\cite{chen2024damsurvey,fang2024dam}.
Automated semantic segmentation of these defects from high-resolution imagery therefore promises substantial gains in inspection efficiency and repeatability.

Recent deep-learning methods have advanced pixel-level defect segmentation considerably.
Encoder--decoder architectures such as U-Net~\cite{ronneberger2015unet} and hierarchical transformers like SegFormer~\cite{xie2021segformer} achieve competitive Intersection-over-Union (IoU) on benchmark crack datasets~\cite{zou2019deepcrack,liu2019deepcrack}.
However, IoU and related region-overlap metrics are insensitive to topological errors: a prediction that fragments a continuous crack into disconnected segments may still score favourably on IoU while misrepresenting the structural extent of the defect~\cite{pantoja2024cracktopology}.
For structural safety assessment, the \emph{connectivity} of a predicted crack is as important as its pixel coverage---a fragmented prediction can lead an engineer to underestimate propagation length and thus misjudge residual capacity.

Topology-preserving losses such as clDice~\cite{shit2021cldice} and persistent-homology penalties~\cite{hu2019topology,clough2020topological} partially address this gap, yet they either incur prohibitive GPU memory for high-resolution dam images or operate on binary masks without multi-class support.
The recently proposed Skeleton Recall Loss (SRL)~\cite{shit2024skeleton} overcomes both limitations by computing recall on CPU-extracted ground-truth skeletons, achieving strong connectivity preservation with over 90\% reduction in computational overhead relative to clDice.
Concurrently, Dynamic Snake Convolution (DSConv)~\cite{qi2023dynamic} reshapes kernel sampling grids to follow elongated, tortuous paths---a morphological prior that aligns naturally with the geometry of cracks.

In this paper we propose \textbf{DSCFormer}, a dual-branch architecture that augments a SegFormer-B2 backbone with a parallel DSConv feature stream and trains with a composite loss incorporating SRL.
We evaluate on a new \textbf{DamSegment} benchmark of 1\,500 annotated images spanning cracks and spalling on concrete dam surfaces.
Our contributions are:
\begin{itemize}
    \item A topology-aware segmentation framework combining dynamic snake convolution with skeleton-based supervision for fine-grained dam defect segmentation, achieving 72.3\% mIoU$_{fg}$ with test-time augmentation.
    \item Integration of Skeleton Recall Loss into a multi-class setting, demonstrating that explicit skeleton supervision yields +4.8 clDice and +5.3 Boundary F1 (BF1) over the SegFormer-B2 baseline without increasing inference cost.
    \item A systematic comparison with foundation model adaptation approaches (SAM-LoRA, DINOv2-LoRA), revealing that task-specific architectural design outperforms parameter-efficient fine-tuning of large vision models on small-scale industrial datasets.
    \item The DamSegment dataset with pixel-level crack and spalling annotations, providing the first public benchmark specifically targeting concrete dam surface defects.
\end{itemize}

\section{Related Work}
\label{sec:related}

\subsection{Concrete Defect Segmentation}
\label{sec:related:defect}

Early CNN-based crack detectors such as CrackNet treated defect identification as patch classification~\cite{zou2019deepcrack}.
Fully convolutional designs subsequently enabled dense pixel-wise prediction: DeepCrack~\cite{liu2019deepcrack} introduced deeply-supervised side outputs fused across encoder--decoder stages, while Zou~\emph{et~al.}~\cite{zou2019deepcrack} exploited hierarchical SegNet features with multi-scale fusion.
These methods established the encoder--decoder paradigm that remains dominant in structural defect segmentation.

More recently, transformer-based architectures have demonstrated improved global context modelling.
CrackFormer~\cite{liu2021crackformer} embeds self-attention within a SegNet-like structure to capture long-range crack continuity, and SegFormer~\cite{xie2021segformer} provides a lightweight hierarchical transformer encoder paired with an all-MLP decoder that achieves strong accuracy--efficiency trade-offs.
Universal architectures such as Mask2Former~\cite{cheng2022mask2former} and hierarchical vision transformers like Swin Transformer~\cite{liu2021swin} further push segmentation accuracy on general benchmarks but have seen limited adoption for infrastructure defects due to annotation scarcity and domain shift.

Applied to dam inspection specifically, Arafin~\emph{et~al.}~\cite{arafin2024concrete} performed multi-class concrete defect segmentation using DeepLabV3+, while Fang and Liu~\cite{fang2024dam} proposed a lightweight MobileNetV2-based architecture for microscopic dam crack quantification.
Chen~\emph{et~al.}~\cite{chen2024damsurvey} surveyed visual perception methods for intelligent dam hub inspection, identifying limited dataset scale and poor topology preservation as key open challenges.

Despite progress, these methods optimise pixel-level overlap metrics (IoU, Dice) and do not explicitly penalise topological disconnections---a critical shortcoming for thin, elongated cracks whose connectivity carries diagnostic significance.

\subsection{Topology-Aware Segmentation}
\label{sec:related:topology}

Topology-aware approaches seek to enforce structural constraints beyond per-pixel accuracy.
Hu~\emph{et~al.}~\cite{hu2019topology} first introduced a differentiable topological loss based on persistent homology, penalising Betti-number discrepancies between predicted and ground-truth segmentations.
Clough~\emph{et~al.}~\cite{clough2020topological} extended persistent-homology supervision to cardiac imaging, demonstrating that topological correctness and pixel accuracy are complementary objectives.
However, persistent-homology computation scales poorly to high-resolution images common in infrastructure inspection.

For tubular structures, Shit~\emph{et~al.}~\cite{shit2021cldice} proposed clDice, a differentiable soft-skeleton metric that guarantees topology preservation up to homotopy equivalence for binary masks.
While elegant, clDice requires GPU-based iterative morphological skeletonisation during training, incurring substantial memory overhead at full resolution and lacking native multi-class support.
The Skeleton Recall Loss~\cite{shit2024skeleton} addresses both limitations: it pre-computes ground-truth skeletons on CPU and formulates a recall objective over skeleton voxels, achieving superior connectivity metrics on five tubular-structure benchmarks with over 90\% reduction in GPU memory relative to clDice.

Recent work on crack-specific topology includes the Crack Topology Score (CTS)~\cite{pantoja2024cracktopology}, a skeleton-matching evaluation protocol weighted by crack length, and boundary-guidance methods~\cite{wang2024dsctopology} that incorporate edge priors into encoder--decoder training.
These advances highlight a growing consensus that pixel-level metrics alone are insufficient for safety-critical crack assessment.

Our work builds upon SRL~\cite{shit2024skeleton} by extending it to the multi-class dam defect setting and pairing it with a geometry-aware feature extractor, yielding complementary gains in both connectivity and boundary accuracy.

\subsection{Deformable and Dynamic Convolutions for Elongated Structures}
\label{sec:related:deformable}

Standard convolutions sample on fixed rectangular grids, an inductive bias poorly matched to thin, curvilinear structures.
Deformable Convolutional Networks (DCN)~\cite{dai2017deformable} introduced learnable offsets that adaptively shift sampling locations, enabling feature extraction along object boundaries.
DCNv2~\cite{zhu2019deformable} augmented offsets with per-sample modulation amplitudes, improving the network's ability to suppress irrelevant context in cluttered scenes.

While DCN offsets are unconstrained and learned purely from data, the Dynamic Snake Convolution (DSConv)~\cite{qi2023dynamic} imposes a \emph{topological geometric constraint}: offsets are parameterised along a one-dimensional path, forcing the kernel to deform into a snake-like shape that naturally traces elongated, tortuous structures such as vessels and cracks.
DSConv demonstrated state-of-the-art vessel segmentation by combining this morphological prior with multi-view feature fusion.

Pyramid pooling~\cite{zhao2017pspnet} and multi-scale attention mechanisms~\cite{liu2021swin} complement deformable operators by providing broader context, yet they do not explicitly encode the elongated geometry of thin defects.
The concurrent DSCformer architecture applies enhanced DSConv to pavement crack segmentation, confirming the effectiveness of snake-shaped receptive fields for this domain.

Our DSCFormer differs from prior DSConv applications in two respects: (i)~we integrate the DSConv branch as a \emph{parallel} stream alongside a SegFormer-B2 encoder rather than replacing standard convolutions wholesale, preserving the transformer's global context; and (ii)~we couple the geometric inductive bias of DSConv with the topological supervision of Skeleton Recall Loss, yielding synergistic improvements in both boundary fidelity and connectivity.


% ===========================================================================
\section{Method}
\label{sec:method}

\subsection{Overview}

Figure~\ref{fig:architecture} illustrates the DSCFormer architecture. Our design philosophy is to enhance topology preservation at two complementary levels: (1)~at the \emph{feature extraction} level through geometry-aware convolutions (DSConv branch), and (2)~at the \emph{training objective} level through explicit skeleton supervision (SRL). Both additions are lightweight and do not increase inference cost for the base segmentation path.

\begin{figure*}[t]
    \centering
    \includegraphics[width=0.95\textwidth]{figures/fig_architecture.png}
    \caption{Overview of the proposed DSCFormer architecture. A SegFormer-B2 encoder extracts hierarchical features, which are fused by an MLP decoder. A parallel Dynamic Snake Convolution branch consumes early encoder stages (1--2) and additively enhances crack channel logits. The model is trained with a composite loss including Skeleton Recall Loss (SRL) activated after epoch~60.}
    \label{fig:architecture}
\end{figure*}

\subsection{SegFormer-B2 Backbone}

We adopt SegFormer-B2~\cite{xie2021segformer} as our backbone, which provides a hierarchical transformer encoder with four stages producing feature maps at resolutions $H/4$, $H/8$, $H/16$, and $H/32$ with channel dimensions 64, 128, 320, and 512, respectively. The lightweight all-MLP decoder fuses multi-scale features through linear projection and upsampling, yielding a strong accuracy--efficiency trade-off. We initialize from weights pretrained on ADE20K at $512{\times}512$ resolution.

\subsection{Dynamic Snake Convolution Branch}
\label{sec:dsconv}

Standard convolutions sample on fixed rectangular grids---an inductive bias poorly matched to thin, curvilinear crack structures. Inspired by Dynamic Snake Convolution~\cite{qi2023dynamic}, we introduce a parallel \emph{CrackSnakeBranch} that performs deformable sampling along the elongated geometry of cracks.

\paragraph{Architecture.} The branch consumes encoder stages 1 and 2 (at $H/4$ and $H/8$ resolution, where crack details are best preserved). Two parallel DSConv paths deform sampling along the $x$-axis and $y$-axis respectively, each with kernel size $K{=}9$ and hidden dimension $d{=}64$. The outputs are summed and projected to a single-channel crack enhancement map $\mathbf{E}_{crack} \in \mathbb{R}^{H \times W}$.

\paragraph{Additive Enhancement.} The enhancement map is added \emph{only} to the crack channel logits of the main segmentation head:
\begin{equation}
    \hat{\mathbf{y}}_{crack} = \mathbf{z}_{crack} + \mathbf{E}_{crack}
\end{equation}
where $\mathbf{z}_{crack}$ denotes the decoder's crack logits. Spalling and background logits remain unmodified, ensuring the branch focuses exclusively on crack topology.

\paragraph{Zero-Initialization.} The output projection layer is initialized with zero weights and bias, so the branch starts as a no-op ($\mathbf{E}_{crack} = \mathbf{0}$) and gradually learns crack-specific deformations during training. This prevents early disruption of the pretrained backbone's predictions.

\subsection{Skeleton Recall Loss}
\label{sec:srl}

To directly enforce topological connectivity, we adopt Skeleton Recall Loss~\cite{shit2024skeleton} adapted to our multi-class setting. Let $S = \{p_1, \ldots, p_{|S|}\}$ denote the set of ground-truth skeleton pixels for the crack class, precomputed offline via morphological skeletonization. SRL penalizes low predicted probability at these critical centerline locations:
\begin{equation}
    \mathcal{L}_{SRL} = 1 - \frac{1}{|S|} \sum_{p \in S} \hat{y}_{crack}(p)
    \label{eq:srl}
\end{equation}
where $\hat{y}_{crack}(p) \in [0,1]$ is the softmax probability for the crack class at pixel $p$.

\paragraph{Advantages over soft-clDice.} Unlike clDice~\cite{shit2021cldice} which requires GPU-based differentiable skeletonization during training, SRL operates on CPU-precomputed skeletons with no additional runtime cost. This enables application to full-resolution $512{\times}512$ images without memory constraints.

\paragraph{Delayed Start.} We activate SRL only after epoch~60 ($\lambda_{SRL} = 0$ for $t < 60$). In early training, the model has not yet learned basic crack features, and skeleton supervision on random predictions causes gradient instability. The delayed start allows stable convergence of the base model before introducing topology regularization.

\subsection{Composite Loss Function}

The full training objective combines cross-entropy for pixel classification, Dice loss for class imbalance mitigation, and SRL for topology preservation:
\begin{equation}
    \mathcal{L} = \lambda_1 \mathcal{L}_{CE} + \lambda_2 \mathcal{L}_{Dice} + \lambda_3 \mathcal{L}_{SRL}
    \label{eq:loss}
\end{equation}
with $\lambda_1 = 0.5$, $\lambda_2 = 0.5$, and $\lambda_3 = 0.05$ (active only for $t \geq 60$). The small SRL weight ensures it acts as a gentle topology regularizer without dominating the gradient.

% ===========================================================================
\section{Experiments}
\label{sec:experiments}

\subsection{Dataset: DamSegment}

We introduce the \textbf{DamSegment} benchmark comprising 1,500 concrete dam surface images at $512{\times}512$ resolution with pixel-level annotations for three classes: background, crack, and spalling. Images are categorized into three difficulty tiers (Easy, Medium, Hard; 500 each) based on defect complexity, lighting conditions, and defect overlap. The dataset is split into train/val/test sets of 1200/150/150 images, stratified by tier and spalling presence to ensure balanced evaluation. Figure~\ref{fig:dataset} shows representative samples.

\begin{figure}[t]
    \centering
    \includegraphics[width=\columnwidth]{figures/fig1_dataset.png}
    \caption{DamSegment dataset samples across difficulty tiers. Left to right: Easy (clear, isolated defects), Medium (moderate complexity), Hard (overlapping defects, poor lighting).}
    \label{fig:dataset}
\end{figure}

\subsection{Implementation Details}

All models are trained on a single NVIDIA GPU using PyTorch. Our DSCFormer uses SegFormer-B2 pretrained on ADE20K-512 as the backbone. Training proceeds for 100 epochs with AdamW optimizer (learning rate $6{\times}10^{-5}$, weight decay $10^{-4}$) and 5-epoch linear warmup. Batch size is 4. The DSConv branch uses $d{=}64$ hidden channels and kernel size $K{=}9$. SRL is activated at epoch 60 with weight 0.05. All baselines are trained with identical data augmentation and evaluation protocols for fair comparison.

\subsection{Evaluation Metrics}

We report six complementary metrics:
\begin{itemize}[nosep]
    \item \textbf{mIoU$_{fg}$}: Mean IoU of crack and spalling classes
    \item \textbf{IoU$_{cr}$}, \textbf{IoU$_{sp}$}: Per-class IoU
    \item \textbf{BF1$_{fg}$}: Boundary F1 score (2px tolerance)
    \item \textbf{clDice$_{fg}$}: Centreline Dice via skeletonization
    \item \textbf{ConnR$_{fg}$}: Connectivity Ratio ($\geq$50\% overlap)
\end{itemize}
BF1 measures edge precision for thin structures, clDice captures topological correctness, and ConnR evaluates connected component preservation---all critical for structural assessment where crack continuity determines diagnostic conclusions.

\subsection{Main Comparison}

Table~\ref{tab:main} presents quantitative results on the DamSegment test set. Our DSCFormer achieves the best performance across all topology-sensitive metrics while maintaining competitive mIoU.

\begin{table*}[t]
\centering
\caption{Quantitative comparison on the DamSegment test set. Best results in \textbf{bold}, second best \underline{underlined}. $\dagger$: parameter-efficient fine-tuning (LoRA, ${\sim}5\%$ trainable params). $\ddagger$: multi-scale + flip test-time augmentation.}
\label{tab:main}
\begin{tabular}{lccccccc}
\toprule
Method & Backbone & mIoU$_{fg}$ & IoU$_{cr}$ & IoU$_{sp}$ & BF1$_{fg}$ & clDice$_{fg}$ & ConnR$_{fg}$ \\
\midrule
U-Net & ResNet-34 & 59.6 & 51.3 & 67.9 & -- & -- & -- \\
DeepLabV3+ & ResNet-50 & 63.4 & 53.3 & 73.5 & 59.2 & -- & -- \\
Mask2Former & Swin-T & 67.9 & 57.1 & 78.6 & 72.5 & \underline{86.6} & 72.6 \\
SegFormer-B2 & MiT-B2 & 70.4 & 56.9 & 83.8 & 66.8 & 81.4 & 80.9 \\
SAM-LoRA$^\dagger$ & ViT-B (SAM) & 68.1 & 55.9 & 80.2 & 68.1 & 82.0 & 81.2 \\
DINOv2-LoRA$^\dagger$ & ViT-B (DINOv2) & 66.7 & 49.1 & 84.2 & 62.4 & 81.7 & 76.5 \\
\midrule
\textbf{Ours (DSCFormer)} & MiT-B2+DSConv & \underline{71.4} & \underline{57.6} & \underline{85.2} & \underline{72.1} & 86.2 & \textbf{82.3} \\
\textbf{Ours + TTA}$^\ddagger$ & MiT-B2+DSConv & \textbf{72.3} & \textbf{59.0} & \textbf{85.6} & \textbf{72.9} & \textbf{86.4} & \underline{79.8} \\
\bottomrule
\end{tabular}
\end{table*}

Compared to the SegFormer-B2 baseline, DSCFormer improves BF1 by \textbf{+5.3} and clDice by \textbf{+4.8}, with an mIoU gain of +1.0 (+1.9 with TTA). We argue that BF1 and clDice improvements are disproportionately valuable for dam inspection: a prediction that correctly identifies crack extent and connectivity enables more reliable structural assessment than one with higher pixel overlap but fragmented topology.

Notably, Mask2Former achieves comparable clDice (86.6) but with substantially lower mIoU (67.9) and ConnR (72.6), suggesting it captures crack shape but misses spalling regions. Foundation models with parameter-efficient fine-tuning (SAM-LoRA, DINOv2-LoRA) also underperform our method: despite leveraging powerful pretrained representations, their generic decoders trained from scratch on only 1,200 images cannot match a fully fine-tuned task-specific architecture. SAM-LoRA achieves 68.1\% mIoU$_{fg}$, 3.3 points below our method, while DINOv2-LoRA reaches only 66.7\%. This demonstrates that for small-scale industrial inspection tasks, domain-specific architectural design (DSConv for cracks) outperforms generic foundation model adaptation.

\subsection{Per-Tier Analysis}

\begin{table}[t]
\centering
\caption{Per-difficulty-tier comparison of SegFormer-B2 baseline versus our DSCFormer. Largest gains appear on Medium and Hard tiers.}
\label{tab:tier}
\begin{tabular}{llccc}
\toprule
Tier & Method & mIoU$_{fg}$ & BF1$_{fg}$ & clDice$_{fg}$ \\
\midrule
\multirow{2}{*}{Easy} & SegFormer-B2 & 74.8 & 65.1 & 84.9 \\
 & \textbf{Ours} & \textbf{76.0} & \textbf{68.3} & \textbf{87.0} \\
\midrule
\multirow{2}{*}{Medium} & SegFormer-B2 & 71.1 & 66.6 & 77.6 \\
 & \textbf{Ours} & \textbf{72.7} & \textbf{77.9} & \textbf{89.6} \\
\midrule
\multirow{2}{*}{Hard} & SegFormer-B2 & 64.4 & 69.4 & 79.4 \\
 & \textbf{Ours} & \textbf{66.1} & \textbf{75.5} & \textbf{84.4} \\
\bottomrule
\end{tabular}
\end{table}

Table~\ref{tab:tier} shows that DSCFormer achieves best results on \emph{all three} difficulty tiers. The largest BF1 improvements occur on Medium (+11.3) and Hard (+6.1) tiers, where crack topology is most challenging due to overlapping defects and poor contrast. This demonstrates that the DSConv branch and SRL are particularly effective when crack geometry is complex and connectivity preservation is most difficult.

\subsection{Ablation Study}

\begin{table}[t]
\centering
\caption{Ablation study showing incremental contributions of each component. DSConv improves crack connectivity; SRL adds consistent gains across all metrics.}
\label{tab:ablation}
\begin{tabular}{lcccccc}
\toprule
Config & mIoU$_{fg}$ & IoU$_{cr}$ & IoU$_{sp}$ & BF1$_{fg}$ & clDice$_{fg}$ & ConnR$_{fg}$ \\
\midrule
(a) SegFormer-B2 & 70.3 & 56.9 & 83.7 & 68.6 & 82.7 & 77.0 \\
(b) + DSConv & 70.4 & 56.9 & 84.0 & 69.5 & 85.5 & \textbf{83.6} \\
(c) + DSConv + SRL & \textbf{71.4} & \textbf{57.6} & \textbf{85.2} & \textbf{72.1} & \textbf{86.2} & 82.3 \\
\bottomrule
\end{tabular}
\end{table}

Table~\ref{tab:ablation} presents the ablation study. Key observations:

\begin{itemize}[nosep]
    \item \textbf{(a)$\to$(b) Adding DSConv}: Minimal mIoU change (+0.1), but +0.9~BF1, +2.8~clDice, and crucially +6.6~ConnR. The deformable snake-shaped sampling grid captures elongated crack geometry that fixed rectangular kernels miss, substantially improving crack connectivity without affecting overall pixel accuracy.
    \item \textbf{(b)$\to$(c) Adding SRL}: +1.0~mIoU, +2.6~BF1, +0.7~clDice. SRL provides consistent gains across \emph{all} metrics by focusing gradient signal on skeleton pixels, improving both segmentation quality and topological completeness. IoU$_{sp}$ also improves (+1.2), indicating SRL's regularizing effect benefits the entire model.
\end{itemize}

The two components are complementary: DSConv improves the \emph{feature representation} (architectural enhancement, primarily ConnR and clDice) while SRL improves the \emph{training signal} (loss enhancement, primarily mIoU and BF1). Their combination yields the best overall results across all six metrics.

\subsection{Qualitative Results}

Figure~\ref{fig:qualitative} presents representative predictions comparing SegFormer-B2 with our DSCFormer. The baseline frequently fragments continuous cracks into disconnected segments, particularly where cracks are thinnest or cross regions of spalling. Our method maintains crack connectivity in these challenging cases, consistent with the quantitative clDice and ConnR improvements.

\begin{figure}[t]
    \centering
    \includegraphics[width=\columnwidth]{figures/fig3_qualitative.png}
    \caption{Qualitative comparison on Hard-tier samples. Top: input image and ground truth. Middle: SegFormer-B2 predictions showing crack fragmentation. Bottom: Our DSCFormer predictions maintaining crack connectivity.}
    \label{fig:qualitative}
\end{figure}

% ===========================================================================
\section{Discussion}
\label{sec:discussion}

\paragraph{Application-Driven Metrics.}
In dam inspection, the primary concern is whether a detected crack represents a continuous propagation path. A topologically correct but spatially imprecise prediction (high clDice, moderate IoU) is more useful for structural assessment than a fragmented prediction with high pixel overlap. Our +5.3~BF1 and +4.8~clDice gains directly translate to more reliable defect characterization.

\paragraph{DSConv Mechanism.}
The DSConv branch's effectiveness stems from its geometric inductive bias: by constraining kernel deformations to follow elongated paths, it captures crack features that rectangular convolutions systematically miss. The zero-initialization ensures the branch only activates when beneficial, avoiding negative transfer for non-crack regions.

\paragraph{SRL vs.\ Soft-clDice.}
While soft-clDice~\cite{shit2021cldice} requires GPU-intensive differentiable skeletonization, our adoption of SRL~\cite{shit2024skeleton} uses precomputed GT skeletons with negligible runtime cost. The delayed-start strategy (epoch 60) is critical: early activation causes training instability as the model attempts topology optimization before learning basic features. This staged approach is simpler and more stable than curriculum-based alternatives we explored.

\paragraph{Foundation Models vs.\ Task-Specific Design.}
Our experiments with SAM-LoRA and DINOv2-LoRA reveal a notable finding: despite leveraging powerful pretrained representations from models trained on hundreds of millions of images, parameter-efficient fine-tuning with LoRA adapters (${\sim}5\%$ trainable parameters) consistently underperforms our fully fine-tuned DSCFormer. SAM-LoRA achieves 68.1\% mIoU$_{fg}$ (3.3 below ours), while DINOv2-LoRA reaches only 66.7\% (4.7 below). We attribute this to two factors: (1)~the FPN decoders trained from scratch on 1,200 images lack sufficient data to fully exploit the frozen encoder features, and (2)~the generic ViT patch embeddings cannot match the crack-specific inductive bias of DSConv. Interestingly, SAM-LoRA+SRL achieves the highest ConnR$_{cr}$ among all methods (86.5\%), suggesting that foundation model features excel at preserving global structure even when local segmentation is imprecise. This finding has practical implications: for small-scale industrial inspection datasets, investing in task-specific architectural design yields better returns than adapting large foundation models.

\paragraph{Limitations.}
The mIoU improvement (+1.0, +1.9 with TTA) is modest compared to topology gains, reflecting that DSConv and SRL primarily address boundary and connectivity rather than region coverage. Crack IoU remains a bottleneck (${\sim}57.6\%$) due to extreme class imbalance (crack pixels constitute ${<}5\%$ of image area). The DSConv branch is currently crack-specific; extending to other elongated defect types (e.g., rebar corrosion) is future work. Additionally, evaluation on other dam inspection datasets would strengthen generalizability claims.

% ===========================================================================
\section{Conclusion}
\label{sec:conclusion}

We presented DSCFormer, a topology-aware segmentation framework for concrete dam defect analysis that achieves state-of-the-art results on our DamSegment benchmark. By combining Dynamic Snake Convolution for geometry-aware feature extraction with Skeleton Recall Loss for explicit connectivity supervision, our method achieves 72.3\% mIoU$_{fg}$ (with TTA), improving BF1 by +5.3 and clDice by +4.8 over the SegFormer-B2 baseline. The two components provide complementary gains---DSConv at the architectural level (ConnR~+6.6) and SRL at the training objective level (mIoU~+1.0)---and both are lightweight additions that do not increase inference cost. Our systematic comparison further demonstrates that task-specific design outperforms foundation model adaptation (SAM-LoRA, DINOv2-LoRA) on small-scale industrial datasets, providing practical guidance for deploying deep learning in data-scarce inspection scenarios.

% ===========================================================================
{\small
\bibliographystyle{ieee_fullname}
\bibliography{refs}
}

\end{document}
